\chapter{Background and Key Concepts}
\label{ch:background}

This chapter gives only the concepts needed to read the method and the results,
and positions the project against prior work.

\section{Two Families of Architectures}

Spoken dialogue systems fall into two families. In the \textbf{cascade} family,
speech recognition, a text dialogue engine and text-to-speech run separately.
This architecture is auditable and replaceable, but component latencies add up
and each stage's errors propagate to the next. In the \textbf{end-to-end}
family, the model goes directly from speech to speech. Natively full-duplex
models such as Moshi model two audio streams in parallel and need no explicit
speech segmentation for turn-taking.

The project definition is end-to-end in spirit, but the bachelor's budget and
the lack of an official Persian trainer made the cascade unavoidable as the
delivered system.

\section{Language Models, Low-Rank Adaptation and Exposure Bias}

The project's text responder is an instruction-following model. Full fine-tuning
of a multi-billion-parameter model on limited hardware was not practical, so the
direct track used LoRA~\cite{hu2022lora}: the weight update is the product of
two low-rank matrices added to the frozen weight.

The decisive point for interpreting the direct-track results is the difference
between training and inference. In training, the model sees the previous
reference token at every step; this is \textbf{teacher forcing}, and the
reported loss is computed under it. At inference, the model must condition on
its own generated tokens. If the two distributions drift apart, errors
accumulate. This mismatch is called \textbf{exposure
bias}~\cite{ranzato2016scheduledsampling}, and its standard remedy is scheduled
sampling~\cite{bengio2015scheduledsampling}. The practical consequence is
explicit: lower teacher-forced loss does not guarantee valid autoregressive
generation.

\section{Recognition, Synthesis and Interruption Features}

The system's speech recognizer is a hybrid RNNT-CTC model of the Conformer
family, run locally with the NeMo stack~\cite{kuchaiev2019nemo}. Text-to-speech
uses the Piper engine~\cite{piper2023} with a Persian voice.

The interruption detector computes classical features on a rolling window:
energy, zero-crossing rate, MFCCs~\cite{davis1980mfcc} and the fundamental
frequency by YIN~\cite{decheveigne2002yin}. The classifier is gradient
boosting~\cite{friedman2001gbm}. Confidence intervals for accuracy are reported
with the Wilson method~\cite{wilson1927interval} and block-wise at the session
level. In human conversation, turn-taking gaps are often around
200 ms~\cite{skantze2021turntaking}, so the interruption decision must be made
much faster than a complete reply is generated.

\section{Position Among Prior Work}
\label{sec:positioning}

Mini-Omni2 and LLaMA-Omni2 have inference repositories but provide no complete,
official speech-to-speech trainer. Freeze-Omni is full duplex, but its trainer
is incomplete for a new language. Qwen2.5-Omni does not list Persian among its
spoken input and output languages, and its 7-billion-parameter version is
outside the target memory range. PersonaPlex~\cite{roy2026personaplex},
built on Moshi, conditions role and voice, but at the time of this project's
decision no official, available trainer for a new language existed. Among the
models reviewed, only the Moshi family met all three conditions of the direct
track: native full duplex, open weights, and an official low-rank adaptation
trainer.

For Persian, single-speaker read-speech resources are no substitute for an
interactive conversational corpus. Table~\ref{tab:positioning} summarizes the
comparison.

\begin{table}[ht]
\centering
\caption{This project compared with prior work on the capabilities of the project definition}
\label{tab:positioning}
\small
\begin{tabular}{L{2.9cm}ccccc}
\toprule
\textbf{Work} & \textbf{Full} & \textbf{Open} & \textbf{Official} &
\textbf{Persian} & \textbf{Memory} \\
              & \textbf{duplex} & \textbf{weights} & \textbf{trainer} &
\textbf{support} & \textbf{range} \\
\midrule
Mini-Omni2   & command-based & yes & no & no & fits \\
LLaMA-Omni2  & no & yes & no & no & fits \\
Freeze-Omni  & yes & yes & partial & no & fits \\
Qwen2.5-Omni & yes & yes & no & no & does not fit \\
Moshi        & native & yes & yes & no & does not fit \\
PersonaPlex  & native & yes & partial & no & does not fit \\
\midrule
This project, cascade & external control & yes & not needed & yes & fits \\
This project, direct  & native & yes & yes & failed & does not fit \\
\bottomrule
\end{tabular}

\vspace{0.5em}
{\footnotesize\itshape The ``memory range'' column is the fit to the 12--24 GB of
the project definition, not the scientific quality of a model. ``Failed'' for
the direct track means failing the autoregressive runtime condition, not a loss
that did not decrease.\par}
\end{table}

\section{Measures and Model Judging}

The mechanical quality of the chain is measured by whether a row passes the
conditions of a Persian transcript, a Persian reply and non-silent audio.
Content preservation of the generated speech is approximated by
re-recognition and character and word error rates; this measure is not human
intelligibility. The semantic quality of replies was measured by an LLM judge
on relevance and coherence, double-blind and with two repetitions. Because the
judge and the candidate are of the same family, the result stays in the class
of ``same-family automatic approximation'' and is no substitute for human
evaluation or the ITU-T P.800 standard~\cite{itu1996p800}. Training,
validation and test splits are always based on the source session identifier,
so that no group leakage occurs.
